\chapter{Team Responsibilities and Use of AI Tools}\label{ch:team}
% Owner: All. Draft 1, 12 Nov 2026; sources: blueprint §15, docs/PLAN.md, AGENTS.md, git history.

\section{Team responsibilities}
Group 15 split the work by layer so that each member owns a part end to end --- code, tests,
report chapters and a demo segment --- while every member reviewed every phase before it was
merged, so that each can explain the whole system.

\begin{longtable}{@{}>{\raggedright\arraybackslash}p{3cm}>{\raggedright\arraybackslash}p{5.6cm}>{\raggedright\arraybackslash}p{3.4cm}>{\raggedright\arraybackslash}p{3.2cm}@{}}
\toprule
\textbf{Member} & \textbf{Code} & \textbf{Report} & \textbf{Tests and demo} \\
\midrule
\endhead
M1 Tushar B \newline Domain and lifecycle & \code{domain/}: users and permissions, the
  \code{Incident} aggregate and its State classes, priority, events; the notification
  subsystem and Telegram alerts; the admin configuration screens & 1 Problem analysis;
  4 UML; 10 OOP and patterns & Domain and notification tests; demo: architecture and
  deployment, Google Sign-In, Workflow 1 \\
M2 Siddhant L Pathak \newline Intelligence and SLA & AI advisor and fallback, similarity and
  duplicates, dependency graph and correlation, predictive SLA risk, assignment strategies,
  SLA engine and monitor, escalation chain, historical seed data & 2 Requirements;
  5 Technology stack & Strategy, similarity (precision and recall), SLA and escalation tests;
  demo: Workflow 2, chat reporting, account lockout \\
M3 Shri Krishna P \newline Platform & API and error mapping, authentication including Google
  Sign-In, repositories, ORM and migrations, audit chain, search and analytics with the
  dashboard, Docker, CI and deployment & 3 Architecture; 6 Repository; 7 Database;
  9 Security & Integration and authorisation-matrix tests; demo: Workflow 3 (duplicates,
  correlation, escalation) \\
M4 Ishaan S \newline Frontend and QA & Design system, Employee and Technician portals, live
  updates, accessibility and phone layout, installable PWA, Playwright suites & 8 Testing;
  final PDF; video editing & End-to-end and accessibility tests, test report; demo:
  administrator portal, health map, search and audit \\
\bottomrule
\end{longtable}

\paragraph{How we worked.} The build ran in six phases (foundations, domain core, persistence and
API, frontend MVP, intelligence and integrations, hardening), each on its own branch and merged
by pull request after all four members reviewed it. \code{docs/PLAN.md} lists every work
package with its owner; \code{docs/DECISIONS.md} records the 45 design decisions taken on the
way and why; \code{docs/PROGRESS.md} says what each phase delivered and what it left for a
human to do.

\section{What each member did}
The work packages below are the ones each member owned in \code{docs/PLAN.md}; a package shared
by two members is listed under both.

\subsection*{Tushar B --- domain model and lifecycle}
\begin{itemize}
  \item Built the core of the domain: the injected \code{Clock} and the \code{ITSMError}
    hierarchy; the \code{User} hierarchy (Employee, Technician, Team Lead, Administrator, IT
    Manager) with the role-to-permission map; the \code{Incident} aggregate with its ten
    \code{State} classes and exactly twenty legal transitions; and the priority strategies with
    category floors and audited overrides.
  \item Defined the repository and adapter interfaces (ports) with in-memory implementations, so
    the domain was tested before any database existed.
  \item Wrote the \code{IncidentService} facade and the hash-chained audit trail, the
    notification subsystem (in-app and email with retries), Telegram alerts, chatbot
    reporting, and the administrator's configuration screens.
  \item Wrote the use case, class and sequence diagrams and report Chapters 1, 4 and 10; shared
    the 80\,\% coverage gate on domain and services with Siddhant.
  \item In the video: the opening, the architecture and deployment, Google Sign-In and
    Workflow~1.
\end{itemize}

\subsection*{Siddhant L Pathak --- intelligence and SLA}
\begin{itemize}
  \item Built the SLA engine (targets per priority and tier, pause and resume) and the SLA
    monitor that warns at 75\,\%, records breaches, escalates up the chain and auto-closes after
    72 hours; the escalation chain itself; and the Skill, Load and Smart assignment strategies
    with their explanations.
  \item Built the intelligent features: duplicate detection (TF-IDF and Gemini embeddings, with
    thresholds tuned on 52 hand-labelled pairs), the service dependency graph and root-cause
    hint, the Gemini category advisor with its keyword-rule fallback, the predictive breach
    risk and the suggested fixes from similar resolved incidents.
  \item Wrote the campus seed data (services and their dependencies, teams, skills and 420
    historical incidents) and, with Shri Krishna, the \code{make demo-reset} demo set-up.
  \item Wrote report Chapters 2 and 5.
  \item In the video: Workflow~2, reporting by chat and the account lockout.
\end{itemize}

\subsection*{Shri Krishna P --- platform}
\begin{itemize}
  \item Set up the repository, the backend skeleton, the cross-platform build commands and
    continuous integration on SQLite and PostgreSQL.
  \item Built persistence (the SQLAlchemy models for 18 tables and the Alembic migrations),
    authentication and security (Argon2 passwords, JWT, lockout, Google Sign-In), the API
    routers with the error-to-HTTP mapping, search with CSV export and the analytics behind
    the dashboard.
  \item Built the authorisation-matrix tests and, with Ishaan, the live WebSocket updates; built
    the administrator's operations screens (users, search, SLA violations, audit log).
  \item Researched free hosting and deployed to Render and Neon with production settings.
  \item Drew the architecture and ER diagrams, wrote the README and report Chapters 3, 6, 7
    and 9.
  \item In the video: Workflow~3 (duplicates, correlation and escalation).
\end{itemize}

\subsection*{Ishaan S --- frontend and quality}
\begin{itemize}
  \item Set up the React frontend and its typed API client, and built the design system shared
    by the three portals.
  \item Built the Employee and Technician portals, the assignment panel with its reasons, the
    live updates in the browser, the service health map, voice input, the light and dark
    themes, the phone layout and the installable web app, and the accessibility checks.
  \item Wrote the Playwright suites that run all three workflows end to end, the 24 documented
    test cases and the generated test report; wrote report Chapter 8 and assembled the final PDF.
  \item Wrote the demonstration script and edited the video.
  \item In the video: the administrator portal, health map, search and audit log.
\end{itemize}

\paragraph{Together.} All four of us reviewed every phase's pull request before it was merged,
rehearsed the demonstration, recorded the video and checked the report.

\section{How AI tools were used}
The requirements allow ``available APIs or packages, AI tools or techniques''. We used AI in two
distinct ways and record both plainly.

\subsection*{AI inside the product}
Google Gemini, through the Google Gen AI SDK, suggests an incident's category with a confidence
score, writes a short summary of how similar past incidents were fixed (citing only those
incidents), and produces text embeddings for duplicate detection. Every one of these is advice:
it is shown with its confidence and can be overridden, every AI decision is written to the
audit trail, and each has a deterministic fallback (keyword rules, a plain list of past fixes,
TF-IDF vectors), so the system works fully without an API key.

\subsection*{AI in building the product}
The code, tests, documentation and most of this report were written with Claude Code, an AI
coding agent, working in cloud sessions on one shared account under rules we wrote for it in
\code{AGENTS.md}:
\begin{itemize}
  \item \textbf{We approved the design first.} The blueprint (\code{docs/blueprint.md}),
    agreed before any code was written, fixes the
    architecture, the domain model, the lifecycle, every business rule and number, and the
    nine patterns. The agent had to follow it and ask before deviating; each agreed deviation
    is logged in \code{docs/DECISIONS.md}.
  \item \textbf{We set the rules of engagement.} Test-first domain rules, \code{make verify}
    before every commit, small reviewed commits on a phase branch, never pushing to
    \code{main}, never reading secrets.
  \item \textbf{We reviewed and merged.} Each phase ended with an independent review of the whole
    phase diff, whose findings were fixed before the pull request; all four members reviewed
    each phase's pull request and a member merged it. The git history shows which commits the
    agent authored and which a member made.
  \item \textbf{We supplied what an agent cannot:} the accounts and keys, the hosting, the
    labelled duplicate pairs to check, the dry runs and the recording.
\end{itemize}
What we learned: an agent is fast at producing code that passes the checks it is given, so the
quality of the result depended on the quality of those checks --- the blueprint's exact rules,
the boundary tests, the layer contracts and the phase reviews. Several real bugs (Chapter~8)
were caught by those reviews, not by the first draft of the code.

